\subsection{The numerical bound for the lowering constant}
\label{archimedean-lowering-product}

The degree-normalized lowering constant reduces to the following
explicit product in the original source \cite{lowering-product-source}.
The estimate below proves the numerical part directly from interlacing.
The identity with the original Bessel scalar remains a separate input.

Let $n\ge k+1$, $k\ge1$, and let
$\lambda_0>\cdots>\lambda_{n-1}$ and $u_1>\cdots>u_k$ be real.
Put $C=\{0,\ldots,k-2\}$ and $N=\{k-1,\ldots,n-1\}$.
Suppose $b\ge B\ge2n$, $b\ge1$, $\gamma\ge4$, and
\[
 \begin{gathered}
 |\lambda_a|,|u_i|\le b,\qquad
 u_{a+1}>\lambda_a>u_{a+2}\quad(a\in C),\\
 0\le\lambda_a\le B/8\quad(a\in C),\qquad
 -2B\le\lambda_a\le-B\quad(a\in N),\\
 \lambda_a-\lambda_c\ge\gamma\quad(a<c,\ a,c\in C).
 \end{gathered}
\]
Let $\mu=\lambda+\delta$, where
\[
 \sum_a|\delta_a|\le q,\qquad
 q\le\gamma/8,\quad q\le B/4,\quad q\le b,
 \qquad \mu_a-\mu_c\ge2\quad(a<c).
\]
For each $s$ set
\[
 R_s=\frac{\prod_{i=1}^k|\mu_s-\tfrac12-u_i|}
                 {\prod_{a\ne s}|\mu_s-1-\mu_a|}.
\]
All denominator factors are at least one, by the last gap condition.

\begin{theorem}\label{original-lowering-product-bound}
Under these hypotheses,
\[
 R_s\le
 \begin{cases}
     8e\sqrt{k}\,b^2,&s\in C,\\
     6e\,b^2,&s\in N.
 \end{cases}
\]
In particular $R_s\le b^4$ throughout the original range
$b\ge B\ge10^6n^7$.
\end{theorem}

\begin{proof}
Put $P(y)=\prod_{a\in C}(y-\lambda_a)$ and
$Q(y)=\prod_i(y-u_i)$.  Interlacing gives the partial fraction identity
\begin{equation}\label{lowering-positive-residues}
 \frac{Q(y)}{P(y)}=y-a-\sum_{t\in C}\frac{w_t}{y-\lambda_t},
 \qquad w_t>0,\qquad |a|\le b,\qquad\sum_tw_t\le b^2.
\end{equation}
Here $w_t=-Q(\lambda_t)/P'(\lambda_t)>0$, and
$a=\sum_i u_i-\sum_t\lambda_t$.  To verify the final bounds without
introducing a factor $k$, form the real symmetric matrix with diagonal
entries $(\lambda_t)_{t\in C},a$ and last column
$(\sqrt{w_t})_{t\in C},a$.  Its other off-diagonal entries are zero.
The determinant formula gives characteristic polynomial $Q$.
Its eigenvalues are therefore the $u_i\in[-b,b]$, so its operator norm
is at most $b$.  The last diagonal entry of its square is
$a^2+\sum_tw_t\le b^2$.  The argument includes $k=1$, with no residues.

Write $x=\mu_s-1$ and $y=x+1/2$.  First suppose $s\in C$.
For $t\in C\setminus\{s\}$, summing consecutive gaps between compact parameters gives
$|\lambda_s-\lambda_t|\ge\gamma|s-t|$.  Since
$|y-\lambda_s|\le q+1/2\le\gamma/4$, we have
\[
       \left|\frac{y-\lambda_s}{y-\lambda_t}\right|\le\frac13.
\]
Multiplying \eqref{lowering-positive-residues} by $y-\lambda_s$
and using its residue sum gives
\[
 \frac{|Q(y)|}{\prod_{t\in C\setminus\{s\}}|y-\lambda_t|}
 \le |y-\lambda_s|\,|y-a|+\sum_tw_t
 \le\frac{25}{4}b^2\le8b^2.
\]
This identity remains valid at $y=\lambda_s$ by polynomial
continuation.  For every $t\ne s$ in $C$,
\[
 |x-\mu_t|\ge\gamma|s-t|-q-1
                   \ge\tfrac12\gamma|s-t|,
 \quad
 \frac{|y-\lambda_t|}{|x-\mu_t|}
          \le1+\frac{|\delta_t|+1/2}{|x-\mu_t|}.
\]
The inequality $\log(1+z)\le z$ and the harmonic-sum bound give
\[
 \begin{split}
 \log\prod_{t\in C\setminus\{s\}}
               \frac{|y-\lambda_t|}{|x-\mu_t|}
 &\le\frac{2q}{\gamma}
              +\frac1\gamma\sum_{t\in C\setminus\{s\}}\frac1{|s-t|}\\
 &\le\frac34+\frac12\log k.
 \end{split}
\]
Thus this product is at most $e\sqrt{k}$.  Discarding the remaining
denominator factors, each at least one, proves the asserted bound for $s\in C$.

Now suppose $s\in N$.  Then $y\le-3B/4-1/2$, while every
$\lambda_t$ with $t\in C$ is nonnegative.  Consequently
\[
 \frac{|Q(y)|}{|P(y)|}
 \le |y-a|+\frac{4b^2}{3B}\le6b^2.
\]
Also $|x-\mu_t|\ge B/2$ for each $t\in C$, and hence
\[
 \log\prod_{t\in C}\frac{|y-\lambda_t|}{|x-\mu_t|}
 \le\frac{2q}{B}+\frac{k-1}{B}\le1.
\]
This gives $R_s\le6eb^2$, again including the empty product $k=1$.

Finally $e<3$, $\sqrt k\le n$, and $B\ge10^6n^7$ imply
$8e\sqrt k\le24n\le b^2$ and $6e\le18\le b^2$.
\end{proof}

The original family supplies precisely these hypotheses after its
integer shift, with $q$ bounded by both the compact-parameter reserve and the
negative-parameter reserve.  The argument even permits real shifts;
no integer or half-integer hypothesis has been substituted for a gap.

For completeness, the exact remaining scalar identification is
\begin{equation}\label{lowering-scalar-identity-input}
 |\kappa_s|^2
       =T\frac{\prod_i|\mu_s-\tfrac12-u_i|}
                    {\prod_{t\ne s}|\mu_s-\mu_t|},\qquad
 \frac{d_\mu}{d_{\mu-e_s}}
       =\frac{\prod_{t\ne s}|\mu_s-\mu_t|}
                    {\prod_{t\ne s}|\mu_s-1-\mu_t|}.
\end{equation}
When these identities hold with the original Bessel normalization and
the same Haar measure, cancellation and
Theorem~\ref{original-lowering-product-bound} give
$|\kappa_s|^2d_\mu/d_{\mu-e_s}\le Tb^4$.
The first identity in \eqref{lowering-scalar-identity-input} is not
proved by the numerical estimate.  In the source it uses the original
Capelli cofactor evaluation and the adjoint identity for the normalized
Bessel functionals.  Those normalization and operator identities,
and the original quotient mass bound, remain explicit obligations.
